\chapter[Considerações finais]{Considerações Finais}
%================================================

\begin{objetivos}
	Ao final desta unidade, o estudante deverá ser capaz de reconhecer a Modelagem de Distribuição de Espécies como um campo integrador entre ecologia, biogeografia, estatística, ciência de dados, mudanças climáticas e conservação; compreender as principais decisões metodológicas discutidas ao longo da apostila; avaliar criticamente os limites e potencialidades dos SDMs; e estruturar interpretações ecológicas responsáveis a partir de modelos espaciais.
\end{objetivos}

\section{Síntese geral da apostila}

Esta apostila foi construída com o objetivo de apresentar a Modelagem de Distribuição de Espécies em R como um fluxo científico completo, crítico e reprodutível. Ao longo das unidades, partimos dos fundamentos conceituais do nicho ecológico e avançamos progressivamente para etapas aplicadas envolvendo dados reais de ocorrência, variáveis ambientais, seleção de preditores, algoritmos estatísticos, aprendizado de máquina, modelos presença-background, ensemble, avaliação, projeções futuras, paleoclima, extrapolação, incertezas e conservação.

O exemplo aplicado com \textit{Dinizia excelsa} permitiu demonstrar que um SDM não deve ser entendido apenas como uma técnica para produzir mapas de adequabilidade ambiental. Ao contrário, ele representa um processo analítico que conecta hipóteses ecológicas, dados espaciais, teoria do nicho, decisões metodológicas, avaliação estatística e interpretação biogeográfica. Nesse sentido, a modelagem de distribuição deve ser vista como uma ferramenta de inferência ecológica espacial, e não como uma simples rotina computacional.

\begin{conceito}[title={\faBook\quad Mensagem central da apostila}]
	Um bom SDM não é definido apenas pela beleza do mapa ou pelo valor de uma métrica. Ele depende da coerência entre pergunta ecológica, qualidade dos dados, definição da área acessível, seleção das variáveis, escolha dos algoritmos, avaliação independente, análise de incertezas e interpretação ecológica responsável.
\end{conceito}

\section{Do nicho ecológico ao espaço geográfico}

Um dos principais aprendizados desta apostila é que a distribuição geográfica de uma espécie não é determinada apenas pelo clima. A presença de uma espécie em determinado local resulta da interação entre condições abióticas, interações bióticas, acessibilidade histórica, dispersão, filtros ambientais, dinâmica populacional e história evolutiva. O arcabouço BAM, discutido nas unidades iniciais, reforça que a distribuição observada é apenas uma parte da distribuição potencial e que a área acessível \(M\) possui papel central na calibração de modelos.

A distinção entre espaço ambiental e espaço geográfico é fundamental. O modelo é ajustado a partir de relações no espaço ambiental, mas seus resultados são projetados no espaço geográfico. Essa transição exige cautela, porque padrões espaciais podem refletir tanto mecanismos ecológicos quanto estrutura amostral, autocorrelação, viés de coleta ou disponibilidade ambiental. Portanto, interpretar mapas de adequabilidade requer sempre retornar à pergunta ecológica original.

\section{Dados de ocorrência e responsabilidade analítica}

A qualidade dos dados de ocorrência é um dos pilares de qualquer SDM. Registros provenientes de inventários, herbários, bancos de dados públicos e plataformas como GBIF possuem enorme valor científico, mas também carregam limitações. Erros taxonômicos, coordenadas imprecisas, duplicatas, registros fora da área de estudo e viés espacial de amostragem podem influenciar diretamente os resultados.

Por isso, a limpeza dos dados não deve ser tratada como etapa burocrática. Ela é parte essencial da inferência ecológica. A rarefação espacial, a verificação taxonômica, a remoção de registros inconsistentes e a documentação das decisões tomadas são procedimentos fundamentais para aumentar a confiabilidade dos modelos.

\begin{atencao}[title={\faExclamationTriangle\quad Dados bons antes de modelos complexos}]
	Nenhum algoritmo compensa dados ruins, área acessível mal definida ou variáveis ambientais incoerentes. Em SDM, a qualidade da inferência começa antes do ajuste do modelo.
\end{atencao}

\section{Variáveis ambientais e interpretação ecológica}

As variáveis ambientais são a ponte entre a ecologia da espécie e a estrutura estatística do modelo. Entretanto, uma variável ambiental não é automaticamente um mecanismo ecológico. Uma variável pode apresentar alta importância no modelo porque resume um gradiente ambiental relevante, porque está correlacionada com outra variável causal, ou porque captura indiretamente processos não medidos.

Ao longo da apostila, enfatizamos a importância de selecionar variáveis com base em critérios ecológicos e estatísticos. A multicolinearidade, avaliada por correlação, VIF e PCA, deve ser tratada com cuidado, especialmente em modelos interpretativos. A seleção de variáveis deve buscar equilíbrio entre parcimônia, representatividade ambiental e plausibilidade ecológica.

\section{Algoritmos como diferentes lentes analíticas}

GLM, GAM, Random Forest, BRT e Maxent representam diferentes formas de aprender relações entre ocorrência e ambiente. GLM oferece maior simplicidade e interpretabilidade; GAM permite respostas suaves e não lineares; Random Forest captura padrões complexos e interações; BRT combina aprendizado sequencial e flexibilidade; Maxent é especialmente relevante para dados presença-background e presença-apenas.

Nenhum algoritmo é universalmente superior. Cada um possui pressupostos, vantagens e riscos. Modelos simples podem ser mais interpretáveis, mas menos flexíveis. Modelos complexos podem melhorar o desempenho preditivo, mas aumentar o risco de sobreajuste. Por isso, a comparação entre algoritmos deve considerar métricas, curvas de resposta, plausibilidade ecológica, estabilidade espacial e capacidade de transferência.

\begin{conceito}[title={\faBook\quad Algoritmos são hipóteses operacionais}]
	Cada algoritmo representa uma hipótese operacional diferente sobre como a espécie responde ao ambiente. Compará-los permite avaliar a robustez dos padrões e não apenas escolher automaticamente o maior valor de AUC.
\end{conceito}

\section{Avaliação, limiares e mapas binários}

A avaliação de modelos é uma etapa indispensável, mas nenhuma métrica isolada é suficiente. AUC, TSS, Kappa, sensibilidade, especificidade, acurácia e Boyce Index respondem a perguntas diferentes. O uso combinado dessas métricas permite uma avaliação mais completa, mas ainda assim deve ser interpretado com cautela.

A conversão de mapas contínuos em mapas binários pode ser útil para análises de área adequada, estabilidade, perda, ganho e conservação. No entanto, mapas binários dependem fortemente do limiar escolhido. Por isso, eles devem ser usados como instrumentos analíticos auxiliares, e não como representação definitiva da distribuição da espécie.

\section{Projeções futuras e mudanças climáticas}

As projeções futuras permitem avaliar como a adequabilidade ambiental pode mudar sob diferentes cenários climáticos. Nesta apostila, os cenários SSP foram apresentados como trajetórias alternativas, não como previsões absolutas. Isso é essencial para evitar interpretações determinísticas.

Modelos projetados para o futuro carregam incertezas associadas aos dados atuais, aos algoritmos, aos modelos climáticos, aos cenários de emissão, à resolução espacial e à extrapolação ambiental. Assim, mapas futuros devem ser interpretados como hipóteses condicionais: se a relação espécie-ambiente permanecer aproximadamente válida e se determinado cenário climático ocorrer, então certas regiões poderão ganhar, perder ou manter adequabilidade ambiental.

\section{Paleoclima e história biogeográfica}

A inclusão do paleoclima amplia a interpretação dos SDMs ao conectar a distribuição atual da espécie a condições ambientais passadas. Projeções para o Holoceno Médio, Último Máximo Glacial e Último Interglacial permitem formular hipóteses sobre estabilidade climática, refúgios históricos e persistência ambiental.

Essas análises são especialmente relevantes para espécies arbóreas amazônicas de longa vida, cujas distribuições contemporâneas podem refletir não apenas o clima atual, mas também legados históricos, dispersão limitada e estabilidade ambiental ao longo do tempo. Entretanto, projeções paleoclimáticas não demonstram ocorrência passada. Elas indicam áreas climaticamente compatíveis com o nicho estimado no presente.

\section{Transferência, extrapolação e incerteza}

A transferência espacial e temporal é uma das aplicações mais poderosas e, ao mesmo tempo, mais frágeis dos SDMs. Quando modelos são projetados para climas futuros ou passados, eles podem encontrar combinações ambientais não representadas na calibração. Nesses casos, as predições tornam-se extrapolativas.

Ferramentas como MESS e MOP ajudam a identificar áreas onde as projeções devem ser interpretadas com cautela. A integração de incerteza algorítmica, incerteza entre cenários, extrapolação e incerteza paleoclimática permite construir mapas mais transparentes e úteis para tomada de decisão.

\begin{atencao}[title={\faExclamationTriangle\quad Alta adequabilidade não significa alta confiança}]
	Uma área pode apresentar alta adequabilidade prevista e, ao mesmo tempo, alta incerteza ou forte extrapolação ambiental. Nessas situações, a interpretação deve ser cautelosa.
\end{atencao}

\section{Aplicações à conservação}

Em conservação, SDMs devem ser usados como camadas de evidência, não como respostas finais. A priorização espacial deve combinar adequabilidade atual, estabilidade futura, estabilidade paleoclimática, baixa incerteza, conectividade, proteção territorial, viabilidade ecológica e conhecimento de campo.

Para espécies arbóreas amazônicas como \textit{Dinizia excelsa}, essa abordagem pode apoiar inventários direcionados, planejamento de áreas prioritárias, restauração ecológica, monitoramento de populações, avaliação de lacunas de proteção e estratégias de conservação sob mudanças climáticas. O valor aplicado do SDM aumenta quando seus resultados são integrados a dados de campo, conhecimento local, políticas públicas e planejamento territorial.

\section{Boas práticas para estudos de SDM}

Com base no fluxo desenvolvido nesta apostila, recomenda-se que estudos de SDM sigam alguns princípios fundamentais:

\begin{enumerate}
	\item formular uma pergunta ecológica clara antes da modelagem;
	\item revisar a taxonomia e a qualidade dos registros de ocorrência;
	\item definir explicitamente a área acessível \(M\);
	\item escolher variáveis ambientais com justificativa ecológica;
	\item avaliar e reduzir multicolinearidade quando necessário;
	\item comparar algoritmos com diferentes níveis de complexidade;
	\item usar validação independente ou particionamento adequado dos dados;
	\item interpretar métricas junto com mapas e curvas de resposta;
	\item avaliar extrapolação em projeções espaciais ou temporais;
	\item quantificar incertezas;
	\item documentar todos os scripts e decisões metodológicas;
	\item comunicar limitações de forma transparente.
\end{enumerate}

\section{Limitações e perspectivas}

Apesar do avanço metodológico apresentado, SDMs continuam sendo simplificações da realidade ecológica. Muitos processos importantes podem não ser representados explicitamente, incluindo dispersão, demografia, interações bióticas, perturbações humanas, mudanças no uso da terra, plasticidade fenotípica, adaptação local e dinâmica populacional.

Perspectivas futuras incluem a integração de SDMs com modelos demográficos, dados genômicos, LiDAR, sensoriamento remoto, traços funcionais, inventários florestais permanentes, modelos de conectividade, cenários de desmatamento e planejamento sistemático da conservação. No contexto amazônico, essa integração é especialmente urgente diante das mudanças climáticas, fragmentação florestal e perda de biodiversidade.

\section{Mensagem final}

A Modelagem de Distribuição de Espécies é uma das ferramentas mais importantes da ecologia espacial contemporânea. Quando aplicada de forma crítica, ela permite transformar dados dispersos de ocorrência e ambiente em hipóteses espaciais úteis para ciência, ensino e conservação.

Contudo, a força de um SDM não está apenas no algoritmo, mas na coerência do processo científico. Um modelo bem construído deve ser ecologicamente fundamentado, estatisticamente avaliado, espacialmente explícito, reprodutível, transparente e interpretado com humildade diante das incertezas.

\begin{conceito}[title={\faLeaf\quad Consideração final}]
	Modelar a distribuição de uma espécie é, em última instância, tentar compreender como a vida ocupa o espaço sob as restrições do ambiente, da história e da mudança. Em um mundo em transformação, SDMs bem conduzidos podem ajudar a antecipar riscos, revelar oportunidades de conservação e orientar decisões mais responsáveis para a biodiversidade.
\end{conceito}

%================================================
